\textbf{Keywords:} Quantization, OSR, NEG FB, STF, NTF, SAR, First
Order, SC SD, CT SD, INCR, FOM

\section{ADC state-of-the-art}\label{adc-state-of-the-art}\index{adc state-of-the-art@ADC state-of-the-art}

The performance of an analog-to-digital converter is determined by the
effective number of bits (ENOB), the power consumption, and the maximum
bandwidth. The effective number of bits contain information on the
linearity of the ADC. The power consumption shows how efficient the ADC
is. The maximum bandwidth limits what signals we can sample and
reconstruct in digital domain.

Many years ago, Robert Walden did a study of ADCs, it is worth looking
up the plots.

1999, R. Walden: Analog-to-digital converter survey and analysis
\citeproc{ref-walden99}{{[}1{]}}

There are obvious trends, the faster an ADC is, the less precise the ADC
is ( lower SNDR). There are also fundamental limits, Heisenberg tells us
that a 20-bit 10 GS/s ADC is impossible, according to Walden.

The plot itself is Figure 4 of that paper, and it is worth looking up:
effective bits against sample rate for every converter Walden could
find, with four sloped lines drawn over the cloud for the limits set by
thermal noise, aperture uncertainty, comparator ambiguity and, at the
top, the ``Heisenberg'' bound. The shape of the cloud is the point. It
runs down and to the right, so every decade of sample rate costs
somewhere around a bit, and no converter sits above the lines.

Murmann's survey, two figures further on, plots the same thing with
twenty-five more years of data, so you can see both what has changed and
what has not.

The uncertainty principle states that the precision we can determine
position and the momentum of a particle is
\[\sigma_x \sigma_p \ge \frac{\hbar}{2}\]. There is a similar relation
of energy and time, given by \[\Delta E \Delta t > \frac{h}{2 \pi}\]
where \(\Delta E\) is the difference in energy, and \(\Delta t\) is the
difference in time.

You should take these limits with a grain of salt. The plot assumes 50
Ohm and 1 V full-scale. As a result, the ``Heisenberg'' line that
appears to be unbreakable certainly is breakable. Just change the
voltage to 100 V, and the number of bits can be much higher. Always
check the assumptions.

A more recent survey of ADCs comes from Boris Murmann. He still
maintains a list of the best ADCs from ISSCC and VLSI Symposium.

\href{https://github.com/bmurmann/ADC-survey}{B. Murmann, ADC
Performance Survey 1997-2023}

A common figure of merit for low-to-medium resolution ADCs is the Walden
figure of merit, defined as

\[ FOM_W = \frac{P}{2^{ENOB} f_s}\]

Below 10 fJ/conv.step is good.

Below 1 fJ/conv.step is extreme.

In the plot below you can see the ISSCC and VLSI ADCs.


{
\centering
\aicfig{media/l6_mwald.png}
}

\emph{Figure 1: Murmann ADC survey: Walden figure of merit versus
Nyquist sample rate for ISSCC and VLSI Symposium ADCs, with the
best-in-class envelope}

\subsection{What makes a state-of-the-art
ADC}\label{what-makes-a-state-of-the-art-adc}

People from NTNU have made some of the world's best ADCs

If you ever want to make an ADC, and you want to publish the
measurements, then you must be better than most. A good algorithm for
state-of-the-art ADC design is to first pick a sample rate with low
number of data (blank spaces in the plot above), then read the papers in
the vicinity of the blank space to understand the application, then set
a target FOM which is best in world, then try and find an ADC
architecture that can achieve that FOM.

That's pretty much the algorithm I, and others, have followed to make
state-of-the-art ADCs. A few of the NTNU ADCs are:

{[}1{]} A Compiled 9-bit 20-MS/s 3.5-fJ/conv.step SAR ADC in 28-nm FDSOI
for Bluetooth Low Energy Receivers \citeproc{ref-wulff17}{{[}2{]}}

{[}2{]} \href{https://ieeexplore.ieee.org/document/9056925}{A 68 dB SNDR
Compiled Noise-Shaping SAR ADC With On-Chip CDAC Calibration}


{
\centering
\aicfig{media/our_work.png}
}

\emph{Figure 2: Walden figure of merit versus Nyquist sample rate with
the NTNU compiled SAR ADCs marked as ``Our work'' near the envelope}

In order to publish, there must be something new. Preferably a new
circuit. Below is the circuit from {[}1{]}. It's a standard
successive-approximation register (SAR) analog-to-digital converter.

The differential input signal is sampled on a capacitor array where the
bottom plate is connected to either VSS or VREF. Once the voltage is
sampled, the comparator will decide whether the differential voltage is
larger, or smaller than 0. Depending on the decision, the MSB capacitors
(left-most) in figure will switch the bottom plate in order to
effectively subtract a voltage equivalent to half the VREF voltage.

The comparator makes another decision, and 1/4'th the VREF voltage is
subtracted or added. Then 1/8'th and so on implementing a binary search
to find the input voltage.

The ``bit-cycling'' (binary-search) loop is self-timed, as such, when
the comparator has made a decision, the next cycle starts.

In (b) we can see the enable flip-flop for the next stage. The CK bar is
the sample clock, as such, A is high during sampling. The output of the
comparator (P and N) is low.

As soon as the comparator makes a decision, P or N goes high, A will be
pulled low, if EI is enabled.

In (c) we can see that the bottom plate of the capacitors \(D_{P0}\),
\(D_{P1}\), \(D_{N0}\), and \(D_{N1}\), are controlled by P and N.

In (d) we can see that the bottom plate of the capacitors also used to
set the comparator clock low again (CO), resetting the comparator, and
pulling P and N low, which in (b) enables the next SAR logic state.

How fast the \(D_{XX}\) settle depend on the size of the capacitors, as
such, the comparator clock will be slow for the MSB, and very fast for
the LSB. This was my main circuit contribution in the paper. I think
it's quite clever, because both the VDD and the capacitor corner will
change the settling time. It's important that the capacitor values fully
settle before the next comparator decision, and as a result of the
circuit in (c,d) the delay is automatically adjusted.

For further details see the paper.


{
\centering
\aicfig{media/fig_sar_logic.png}
}

\emph{Figure 3: SAR ADC schematic: (a) capacitor array with self-timed
SAR logic chain and comparator, (b) enable flip-flop, (c) bottom-plate
switching of the CDAC, (d) comparator clock generation}

For state-of-the-art ADC papers it's not sufficient with the idea, and
simulation. There must be proof that it actually works. No-one will
really believe that the ADC works until there is measurements of an
actual taped out IC.

Below you can see the layout of the IC I made for the paper. Notice that
there are 9 ADCs. I had many ideas that I wanted to try out, and I was
not sure what would actually be state of the art. As a result, I taped
out multiple ADCs.


{
\centering
\aicfig{media/l06_fig_layout.png}
}

\emph{Figure 4: Layout of the test chip with nine ADC variants inside
the pad ring}

The two ADCs that I ended up using in the paper are shown below. The one
on the left was made with 180 nm IO transistors, while the one on the
right was made with core-transistors. Notice that the layout of the two
is quite similar.


{
\centering
\aicfig{media/l06_fig_toplevel.png}
}

\emph{Figure 5: Layout of the two compiled SAR ADCs with comparator,
logic, CDAC and switch: (a) 180 nm IO-transistor version (40 x 106 um),
(b) core-transistor version (39 x 80 um)}

Once taped out, and many months of waiting, a few months of measurement
in the lab, I had some results that would be good enough to qualify for
the best conference, and luckily the best journal.


{
\centering
\aicfig{media/l06_core_meas_tikz.png}
}

\emph{Figure 6: Measured performance of the core-transistor ADC: (a, b)
output spectra at 0.69 V and 0.47 V supply, (c) peak ENOB versus VDD,
(d) SNDR and SFDR versus input frequency}

Comparing my ADCs to others, we can see that the FOM is similar to
others. Based on the FOM it might not be clear why the paper was
considered state-of-the-art.

The circuit technique mentioned above would not have been enough to
qualify. The big thing was the ``Compiled'' line. Compared to the other
``Compiled'' mine was 300 times better, and on par with other
state-of-the-art.


{
\centering
\aicfig{media/l06_jssc_table.png}
}

\emph{Figure 7: Comparison table against state-of-the-art SAR ADCs,
where ``This work'' stands out as the only compiled ADC with competitive
figure of merit}

The big thing was how I made the ADC. I started with a definition of a
transistor, as shown below


{
\centering
\aicfig{media/l06_fig_dmos.png}
}

\emph{Figure 8: Transistor definition used by the layout compiler:
diffusion (OD), contacts (CO), poly (PO) and metal 1 (M1) placed on a
vertical and horizontal grid}

And then wrote a compiler (in Perl, later C++
\href{https://github.com/wulffern/ciccreator}{ciccreator}) to compile a
object definition file, a SPICE netlist and a technology rule file into
the full ADC layout.

In (a) you can see one of the cells in the SAR logic, (b) is the spice
file, and (c) is the definition of the routing. The numbers to the right
in the routing creates the paths shown in (d).


{
\centering
\aicfig{media/l06_fig_saremx.png}
}

\emph{Figure 9: Compiled SAR logic cell: (a) 3D view of the layout, (b)
SPICE netlist, (c) object definition with routing rules, (d) the
resulting routed layout}

The implementation is the
\href{https://github.com/wulffern/sun_sar9b_sky130nm/blob/main/cic/ip.spi}{SPICE
netlist}, and the
\href{https://github.com/wulffern/sun_sar9b_sky130nm/blob/main/cic/ip.json}{object
definition file} (JSON) and the
\href{https://github.com/wulffern/sun_sar9b_sky130nm/blob/main/cic/sky130.tech}{rule
file}.

What I really like is the fact that the compilation could generate GDSII
or SKILL, or these days, Xschem schematics and Magic layout.


{
\centering
\aicfig{media/l06_fig_process.png}
}

\emph{Figure 10: Compiled ADC design flow: architecture, implementation
files (SPICE netlist, object definition, technology file), compilation
to GDSII or SKILL, and physical verification}

The cool thing with a compiled ADC is that it's easy to port between
technologies. Since the original ADC, I've ported the ADC to multiple
closed PDKs (22 nm FDSOI, 22 nm, 28 nm, 55 nm, 65 nm and 130nm). In the
summer of 2022 I made an open source port to skywater 130nm.

\href{https://github.com/wulffern/sun_sar9b_sky130nm/}{SUN\_SAR9B\_SKY130NM}


{
\centering
\aicfig{media/l00_SAR9B_CV.png}
}

\emph{Figure 11: The SAR ADC ported to skywater 130 nm: Magic layout of
the ADC core and ngspice transient simulation of a conversion}

One of my Ph.D students built on-top on my work, and made a noise-shaped
compiled SAR ADC, shown below, more on that later.


{
\centering
\aicfig{media/harald_layout.png}
}

\emph{Figure 12: Noise-shaping compiled SAR ADC: die photo of the two
ADC instances and the 116 um x 202 um core layout with CDAC, SAR logic,
loop filter, OTAs and code correction}

\subsection{High resolution FOM}\label{high-resolution-fom}\index{high resolution fom@High resolution FOM}

For high-resolution ADCs, it's more common to use the Schreier figure of
merit, which can also be found in

\href{https://web.stanford.edu/~murmann/adcsurvey.html}{B. Murmann, ADC
Performance Survey 1997-2022 (ISSCC \& VLSI Symposium)}

The Walden figure of merit assumes that thermal noise does not constrain
the power consumption of the ADC, which is usually true for
low-to-medium resolution ADCs. To keep the Walden FOM you can double the
power for a one-bit increase in ENOB. If the ADC is limited by thermal
noise, however, then you must quadruple the capacitance (reduce \(kT/C\)
noise power) for each 1-bit ENOB increase. Accordingly, the power must
also go up four times.

For higher resolution ADC the power consumption is set by thermal noise,
and the Schreier FOM allows for a 4x power consumption increase for each
added bit.

\[FOM_S = SNDR + 10\log\left(\frac{f_s/2}{P}\right)\]

Above 180 dB is extreme


{
\centering
\aicfig{media/l6_msch.png}
}

\emph{Figure 13: Murmann ADC survey: Schreier figure of merit versus
Nyquist sample rate, where the envelope flattens around 185 dB for
thermal-noise limited ADCs}

\section{Quantization}\label{quantization}\index{quantization@Quantization}

Sampling turns continuous time into discrete time. Quantization turns
continuous value into discrete value. Any complete ADC is always a
combination of sampling and quantization.

In our mathematical drawings of quantization we often define \(y[n]\) as
the output, the quantized signal, and \(x[n]\) as the discrete time,
continuous value input, and we add some ``noise'', or ``quantization
noise'' \(e[n]\), where \(x[n] = y[n] - e[n]\).


{
\centering
\aicfig{media/l6_adc_tikz.png}
}

\emph{Figure 14: Linear model of quantization: the quantization noise
e{[}n{]} is added to the input x{[}n{]} to form the output y{[}n{]}}

Maybe you've even heard the phrase ``Quantization noise is white'' or
``Quantization noise is a random Gaussian process''?

I'm here to tell you that you've been lied to. Quantization noise is not
white, nor is it a Gaussian process. Those that have lied to you may say
``yes, sure, but for high number of bits it can be considered white
noise''. I would say that's similar to saying ``when you look at the
earth from the moon, the surface looks pretty smooth without bumps, so
let's say the earth is smooth with no mountains''.

I would claim that it's an unnecessary simplification. It's obvious to
most that the earth would appear smooth from really far away, but they
would not be surprised by Mount Everest, since they know it's not
smooth. An Alien that has been told that the earth is smooth, would be
surprised to see Mount Everest.

But if Quantization noise is not white, what is it?

The figure below shows the input signal x and the quantized signal y.


{
\centering
\aicfig{media/l6_ct_tikz.png}
}

\emph{Figure 15: A continuous time sinusoid input x (blue) and the
quantized output y (red)}

To see the quantization noise, first take a look at the sample and held
version of \(x\) in green in the figure below. The difference between
the green (\(x\) at time n) and the red (\(y\)) would be our
quantization noise \(e\)

The quantization noise is contained between \(+\frac{1}{2}\) Least
Significant Bit (LSB) and \(-\frac{1}{2}\) LSB.

This noise does not look random to me, but I can't see what it is, and
I'm pretty sure I would not be able to work it out either.


{
\centering
\aicfig{media/l6_cten_tikz.png}
}

\emph{Figure 16: Sample-and-held input (green), quantized output (red),
and the resulting quantization error e{[}n{]}, bounded by plus/minus
half an LSB}

Luckily, there are people in this world that love mathematics, and that
can delve into the details and figure out what \(e[n]\) is. A guy called
Blachman wrote a paper back in 1985 on quantization noise.

See The intermodulation and distortion due to quantization of sinusoids
\citeproc{ref-blachman85a}{{[}3{]}} for details

In short, the quantizer \emph{output} for a sinusoidal input is an exact
harmonic series,

\[y(t) = \sum_{p=1}^\infty{A_p\sin{p\omega t}}\]

and the quantization noise is what is left when the input is taken back
out,

\[ e_n(t) = y(t) - A\sin{\omega t} \]

which is the same series with the \(\delta_{p1}A\) term removed. It is
worth keeping the two apart. The series below carries the fundamental at
nearly full amplitude, so it is emphatically not something with zero
mean and variance \(\Delta^2/12\); the \emph{difference} is.

where p is the harmonic index, and

\[
A_p = 
\begin{cases}
\delta_{p1}A  + \sum_{m =
  1}^\infty{\frac{2}{m\pi}J_p(2m\pi A)} &, p = \text{ odd} \\
 0 &, p = \text{ even}
\end{cases}
\]

\[
\delta_{p1}
\begin{cases}
1 &, p=1 \\
0 &, p \neq 1
\end{cases}
\]

and \[J_p(x)\] is a Bessel function of the first kind, \emph{A} is the
amplitude of the input signal.

If we approximate the amplitude of the input signal as

\[A = \frac{2^n - 1}{2} \approx 2^{n-1}\]

where n is the number of bits, we can rewrite as

\[y(t) = \sum_{p=1}^\infty{A_p\sin{p\omega t}}\]

\[ A_p = \delta_{p1}2^{n-1} + \sum_{m=1}^\infty{\frac{2}{m\pi}J_p(2m\pi
  2^{n-1})},  p=odd\]

Obvious, right?

I must admit, it's not obvious to me. But I do understand the
implications. The quantization noise is an infinite sum of input signal
odd harmonics, where the amplitude of the harmonics is determined by a
sum of a
\href{https://en.wikipedia.org/wiki/Bessel_function\#Bessel_functions_of_the_first_kind}{Bessel
function}.

A Bessel function of the first kind looks like this


{
\centering
\aicfig{media/bessel_tikz.png}
}

\emph{Figure 17: Bessel functions of the first kind, J0(x), J1(x) and
J2(x), showing the oscillatory behavior that shapes the quantization
noise harmonics}

So I would expect the amplitude to show signs of oscillatory behavior
for the harmonics. That's the important thing to remember. The
quantization noise is \textbf{odd harmonics of the input signal}

The mean value is zero


{
\centering
\aicfig{media/quant_noise_tikz.png}
}

\emph{Figure 18: What the Bessel series actually says: a sine through a
3-bit quantizer, the error it leaves, and the odd harmonic amplitudes
\(A_p\) from the formula above, against the flat floor the
\(\Delta^2/12\) model would predict}

Here is the formula drawn out. On the left, a sine through a three bit
quantizer, and underneath it the error it leaves behind. The error is
not random: it repeats exactly once per signal period, which is the
whole reason its spectrum can only contain harmonics of the signal.

On the right, the amplitudes \(A_p\) evaluated from the Bessel sum. They
agree with a direct FFT of the quantized sine to about a tenth of a
percent, so this is not an approximation of the noise, it \emph{is} the
noise. Compare the spikes with the dashed line, which is where the
\(\Delta^2/12\) white noise model would put a flat floor: the total
power is the same, but the distribution is nothing like it.

If you want to feel this rather than read it, the
\href{https://wulffern.github.io/aic2026/assets/examples/bessel-quantization.html}{interactive
version} puts the bit count on a slider. Walk it up and watch the
harmonics fall and crowd together, until somewhere around eight bits
calling them a noise floor finally becomes fair.

Watch the rate they fall at, because it is not the 6 dB per bit you
might expect. An individual harmonic drops about \textbf{9 dB per bit},
tending to \(30\log 2 = 9.03\) dB. The 6 dB per bit belongs to the
\emph{total} quantization noise power, and the difference between the
two is the whole reason the noise floor idea works at all: each harmonic
falls 9 dB, but the number of harmonics crammed below \(f_s/2\) doubles,
which adds 3 dB back. Nine down, three up, six net. So as you add bits
the spectrum does not merely shrink, it redistributes --- fewer and
fewer dB in each of more and more spikes, until no individual spike is
worth naming and the aggregate is all that is left.

\[\overline{e_n(t)} = 0 \]

and variance (mean square, since mean is zero), or noise power, can be
approximated as

\[\overline{e_n(t)^2} = \frac{\Delta^2}{12}\]

\subsection{Signal to Quantization noise
ratio}\label{signal-to-quantization-noise-ratio}

Assume we wanted to figure out the resolution, or effective number of
bits for an ADC limited by quantization noise. A power ratio, like
signal-to-quantization noise ratio (SQNR) is one way to represent
resolution.

Take the signal power, and divide by the noise power

\[ SQNR = 10 \log\left(\frac{A^2/2}{\Delta^2/12}\right) = 10 \log\left(\frac{6 A^2}{\Delta^2}\right) \]

\[ \Delta = \frac{2A}{2^B}\]

\[ SQNR = 10 \log\left(\frac{6 A^2}{4 A^2/2^{2B}}\right) = 20 B \log 2 + 10 \log 6/4\]

\[ SQNR  \approx 6.02 B + 1.76\]

You may have seen the last equation before, now you know where it comes
from.

\subsection{Understanding
quantization}\label{understanding-quantization}

Below I've tried to visualize the quantization process
\href{https://github.com/wulffern/aic2026/blob/main/ex/q.py}{q.py},
which also exists as an
\href{https://wulffern.github.io/aic2026/assets/examples/quantization.html}{interactive
page} where the number of bits is a slider and the SQNR is measured for
you.

The left most plot is a sinusoid signal and random Gaussian noise. The
signal is not a continuous time signal, since that's not possible on a
digital computer, but it's an approximation.

The plots are FFTs of a sinusoidal signal combined with noise. These are
complex FFTs, so they show both negative and positive frequencies; the
x-axis is normalized to each record's sample rate, from \(-f_s/2\) to
\(+f_s/2\). When the text below talks about \emph{bins}, multiply
\(f/f_s\) by the record length (2048 after sampling) to get the bin
number. Notice that there are two spikes, which should not be
surprising, since a sinusoidal signal is a combination of two
frequencies.

\[
sin(x) = \frac{e^{ix} - e^{-ix}}{2i}
\]

The second plot from the left is after sampling, notice that the noise
level increases. The rise is exactly \(10\log(4) = 6.02\) dB, and it is
worth being careful about why, because there are two tempting
explanations and they are not two effects to be added together.

Both are named in the same expression. With a Hann window and the peak
normalisation that \texttt{freqDomain} uses, the noise floor per bin
relative to the tone is

\[ \text{floor} = \frac{6\sigma^2}{M} \]

where \(\sigma\) is the time-domain noise and \(M\) the record length.
The record length is right there in the denominator, so a four times
shorter record does raise the floor 6 dB. But \(\sigma^2\) is in the
numerator, and whether it survives decimation is exactly the folding
question. Decimating without an anti-alias filter preserves the total
noise power, \(\sigma\) does not change, and the floor rises the full 6
dB. Put an ideal brick-wall filter in front and \(\sigma^2\) drops by
four as well, the two effects cancel term for term, and the floor does
not move at all.

So the answer is 6.02 dB in total and the two mechanisms must not be
added: they are two ways of booking the same rise, and the experiment
that separates them is to filter before decimating. It is the first
appearance in this chapter of a rule worth keeping --- throwing samples
away only helps if something band-limits the noise first.

The right plot is after quantization, where I've used the function
below.

\begin{Shaded}
\begin{Highlighting}[]
\KeywordTok{def}\NormalTok{ adc(x,bits):}
\NormalTok{    levels }\OperatorTok{=} \DecValTok{2}\OperatorTok{**}\NormalTok{bits}
\NormalTok{    delta }\OperatorTok{=} \DecValTok{2}\OperatorTok{/}\NormalTok{levels}
\NormalTok{    y }\OperatorTok{=}\NormalTok{ np.floor(x}\OperatorTok{/}\NormalTok{delta)}\OperatorTok{*}\NormalTok{delta }\OperatorTok{+}\NormalTok{ delta}\OperatorTok{/}\DecValTok{2}
    \ControlFlowTok{return}\NormalTok{ np.clip(y, }\OperatorTok{{-}}\DecValTok{1} \OperatorTok{+}\NormalTok{ delta}\OperatorTok{/}\DecValTok{2}\NormalTok{, }\DecValTok{1} \OperatorTok{{-}}\NormalTok{ delta}\OperatorTok{/}\DecValTok{2}\NormalTok{)}
\end{Highlighting}
\end{Shaded}

A B-bit converter has exactly \(2^B\) levels, spaced by
\(\Delta = 2/2^B\), sitting half a step off zero. The obvious one-liner
\texttt{np.round(x*2**bits)/2**bits} looks like a quantizer but is not
one: its step is \(2^{-B}\), so at one bit it produces five output
levels over \(\pm 1\) rather than two, and the measured SQNR comes out a
whole bit too good. The plots below are a real one bit converter - two
levels, one comparator.

I really need you to internalize a few things from the right most plot.
Really think through what I'm about to say.

Can you see how the noise (what is not the two spikes) is not white?
White noise would be flat in the frequency domain, but the noise is not
flat.

Notice also that this is a genuine one bit quantizer: two levels, so the
output is a square wave, and what you are looking at is its odd harmonic
series - exactly the \(A_p\) of the Bessel formula above, folded back
into the band wherever a harmonic lands above \(f_s/2\).


{
\centering
\aicfig{media/l6_q_1_tikz.png}
}

\emph{Figure 19: FFT of a sinusoid with noise as continuous value
(left), after sampling (middle), and after 1-bit quantization (right),
where the quantization noise shows up as distinct harmonic spikes rather
than a white noise floor}

If you run the python script you can zoom in and check the highest
spikes. A 1-bit quantizer is a sign detector, so its output for a sine
input is a square wave, and a square wave has only odd harmonics with
amplitudes falling as \(1/p\). That is exactly what the plot shows. The
fundamental is at bin 127, and the measured spikes are

\begin{longtable}[]{@{}llll@{}}
\toprule\noalign{}
harmonic & \(20\log(1/p)\) & measured & bin \\
\midrule\noalign{}
\endhead
\bottomrule\noalign{}
\endlastfoot
3 & \(-9.54\) dB & \(-9.54\) dB & 381 \\
5 & \(-13.98\) dB & \(-13.98\) dB & 635 \\
7 & \(-16.90\) dB & \(-16.90\) dB & 889 \\
9 & \(-19.08\) dB & \(-19.08\) dB & \textbf{905} \\
11 & \(-20.83\) dB & \(-20.83\) dB & \textbf{651} \\
\end{longtable}

Agreement to the second decimal with a formula you can write down from
memory is a good afternoon's work. But look at the last two rows. The
ninth harmonic should be at \(9 \times 127 = 1143\) and the eleventh at
\(11 \times 127 = 1397\), and neither bin exists --- the record is only
2048 points, so anything above 1024 is above half the sample rate. They
fold: \(2048 - 1143 = 905\) and \(2048 - 1397 = 651\). The harmonics
march monotonically down in amplitude, but from bin 889 onwards they
march \emph{backwards} along the frequency axis, which is why a spectrum
plot of an aliased converter can look like nonsense until you work out
which harmonic each spike really is.

This is worth internalising, because in a real converter you do not get
to compare against a formula. A spike at 651 with nothing at 1397 is not
evidence of some mysterious distortion mechanism; it is the eleventh
harmonic, folded. If you change the python script to reduce the
frequency, \texttt{fdivide=2**9}, and increase number of points,
\texttt{N=2**16}, as in the plot below, the eleventh harmonic no longer
folds, and you'll see it directly at bin 1397.


{
\centering
\aicfig{media/l6_q_1_fharm_tikz.png}
}

\emph{Figure 20: The same 1-bit quantization with lower input frequency
and a 16384-point FFT, where the 11th harmonic appears directly at bin
1397 instead of folding}

All the other spikes are the odd harmonics above the sample rate that
fold. The infinite sum of harmonics will fold, some in-phase, some out
of phase, depending on the sign of the Bessel function.

From the function for the amplitude of the quantization noise for
harmonic indices higher than \(p=1\)

\[ A_p =  \sum_{m=1}^\infty{\frac{2}{m\pi}J_p(2m\pi 2^{n-1}  ) }\text{,  p=odd} \]

we can see that the input to the Bessel function increases faster for a
higher number of bits \(n\). As such, from the Bessel function figure
above, I would expect that the sum of the Bessel function is a lower
value. Accordingly, the quantization noise reduces at higher number of
bits.

A consequence is that the quantization noise becomes more and more
uniform, as can be seen from the plot of a 10-bit quantizer below.
That's why people say ``Quantization noise is white'', because for a
high number of bits, it looks white in the FFT.


{
\centering
\aicfig{media/l6_q_10_tikz.png}
}

\emph{Figure 21: FFT of the same signal with a 10-bit quantizer, where
the quantization noise is closer to uniform and looks almost white}

\subsection{Why you should care about quantization
noise}\label{why-you-should-care-about-quantization-noise}

So why should you care whether the quantization noise looks white, or
actually is white? A class of ADCs called oversampling and sigma-delta
modulators rely on the assumption that quantization noise \textbf{is}
white. In other words, the cross-correlation between noise components at
different time points is zero. As such the noise power sums as a sum of
variance, and we can increase the signal-to-noise ratio.

\textbf{We} know that assumption to be wrong though,
\textbf{quantization noise is not white}. For noise components at
harmonic frequencies the cross-correlation will be high. As such, when
\textbf{we} design oversampling or sigma-delta based ADC \textbf{we}
will include some form of dithering (making quantization noise whiter).
For example, before the actual quantizer we inject noise, or we make
sure that the thermal noise is high enough to dither the quantizer.

Everybody that thinks that quantization noise \textbf{is} white will
design non-functioning (or sub-optimal) oversampling and sigma-delta
ADCs. That's why you should care about the details around quantization
noise.

\section{Oversampling}\label{oversampling}\index{oversampling@Oversampling}

Here is the question this section answers. If the quantizer's noise is a
fixed amount that we cannot reduce, can we win anything by sampling
faster than the signal actually demands, and then averaging the extra
samples away?

The answer is yes, and the reason is that signal and noise behave
differently under summation. A slow signal is nearly the same from one
sample to the next, so summing samples adds it up coherently. Noise is
not, so it adds up incoherently. Everything below is bookkeeping on that
one asymmetry.

Assume a signal \(x[n] = a[n] + b[n]\) where \(a\) is a sampled sinusoid
and \(b\) is a random process where cross-correlation is zero for any
time except for \(n=0\). Assume that we sum two (or more) equally spaced
signal components, for example

\[y = x[n] + x[n+1]\]

What would the signal to noise ratio be for \(y\)?

\subsection{Noise power}\label{noise-power}\index{noise power@Noise power}

Our mathematician friends have looked at this, and as long the noise
signal \(b\) \textbf{is random} then the noise power for the oversampled
signal \(b_{osr} = b[n] + b[n+1]\) will be

\[ \overline{b_{osr}^2} = OSR \times \overline{b^2} \]

where OSR is the oversampling ratio. If we sum two time points the
\(OSR=2\), if we sum 4 time points the \(OSR=4\) and so on.

For fun, let's go through the mathematics

Define \(b_1 = b[n]\) and \(b_2 = b[n+1]\) and compute the noise power

\[
\overline{(b_1 + b_2)^2} = \overline{b_1^2 + 2b_1b_2 + b_2^2}
\]

Let's replace the mean with the actual function

\[
\frac{1}{N}\sum_{n=0}^N{\left(b_1^2 + 2b_1b_2 + b_2^2\right)}
\]

which can be split up into

\[
\frac{1}{N}\sum_{n=0}^N{b_1^2} + \frac{1}{N}\sum_{n=0}^N{2b_1b_2} + \frac{1}{N}\sum_{n=0}^N{b_2^2}
\]

we've defined the cross-correlation to be zero, as such

\[
\overline{(b_1 + b_2)^2} = \frac{1}{N}\sum_{n=0}^N{b_1^2} + \frac{1}{N}\sum_{n=0}^N{b_2^2} =  \overline{b_1^2} + \overline{b_2^2}
\]

but the noise power of each of the \(b\)'s must be the same as \(b\), so

\[
\overline{(b_1 + b_2)^2} = 2\overline{b^2}
\]

\subsection{Signal power}\label{signal-power}\index{signal power@Signal power}

For the signal \(a\) we need to calculate the increase in signal power
as OSR increases.

I like to think about it like this. \(a\) is low frequency, as such,
samples \(n\) and \(n+1\) is pretty much the same value. If the sinusoid
has an amplitude of 1, then the amplitude would be 2 if we sum two
samples. As such, the amplitude must increase with the OSR.

The signal power of a sinusoid is \(A^2/2\), accordingly, the signal
power of an oversampled signal must be \((OSR \times A)^2/2\).

\subsection{Signal to Noise Ratio}\label{signal-to-noise-ratio}\index{signal to noise ratio@Signal to Noise Ratio}

Take the signal power to the noise power

\[
\frac{(OSR \times A)^2/2}{OSR \times \overline{b^2}} = OSR \times \frac{A^2/2}{\overline{b^2}}
\]

We can see that the signal to noise ratio increases with increased
oversampling ratio, \textbf{as long as the cross-correlation of the
noise is zero}

\subsection{Signal to Quantization Noise
Ratio}\label{signal-to-quantization-noise-ratio-1}

Now put the two halves together. The quantizer always makes the same
total amount of noise, \(\Delta^2/12\), and always spreads it evenly
from zero to \(f_s/2\). Sampling faster than we need does not reduce
that total, it only spreads it over a wider band, so the part that lands
inside the band we actually care about shrinks by the oversampling
ratio.

in-band quantization noise for an oversampling ratio (OSR)

\[ \overline{e_n(t)^2} =\frac{\Delta^2}{12 OSR}\]

That single division by OSR is the whole of oversampling, and everything
else in this section is arithmetic on it. Dividing the noise power by
OSR multiplies the ratio by OSR:

\[ SQNR = 10 \log\left(\frac{6 A^2}{\Delta^2/OSR}\right) = 10 \log\left(\frac{6 A^2}{\Delta^2}\right) + 10 \log(OSR)\]

\[ SQNR \approx 6.02B + 1.76 + 10 \log(OSR)\]

so the SQNR improves by \(10\log(2) \approx 3\) dB for an OSR of 2, and
\(10\log(4) \approx 6\) dB for an OSR of 4 --- 3 dB for every doubling.

\[ 10 \log(2) \approx 3 dB\]

\[ 10 \log(4) \approx 6 dB\]

Compare that with the 6.02 dB a real bit is worth and the exchange rate
falls out: three decibels is half a bit, so oversampling buys 0.5-bit
per doubling of OSR

It is a poor exchange rate, and it is worth feeling how poor before
moving on. Going from an 8-bit converter to a 12-bit one by oversampling
alone needs \(OSR = 2^8 = 256\), which means running the analog front
end 256 times faster than the signal requires. That is the wall that
noise shaping exists to get around.

\subsection{Python oversample}\label{python-oversample}\index{python oversample@Python oversample}

There are probably more elegant (and faster) ways of implementing
oversampling in python, but I like to write the dumbest code I can,
simply because dumb code is easy to understand.

Below you can see an example of oversampling. The \texttt{oversample}
function takes in a vector and the OSR. For each index it sums OSR
future values.

\begin{Shaded}
\begin{Highlighting}[]
\KeywordTok{def}\NormalTok{ oversample(x,OSR):}
\NormalTok{    N }\OperatorTok{=} \BuiltInTok{len}\NormalTok{(x)}
\NormalTok{    y }\OperatorTok{=}\NormalTok{ np.zeros(N)}

    \ControlFlowTok{for}\NormalTok{ n }\KeywordTok{in} \BuiltInTok{range}\NormalTok{(}\DecValTok{0}\NormalTok{,N):}
        \ControlFlowTok{for}\NormalTok{ k }\KeywordTok{in} \BuiltInTok{range}\NormalTok{(}\DecValTok{0}\NormalTok{,OSR):}
\NormalTok{            m }\OperatorTok{=}\NormalTok{ n}\OperatorTok{+}\NormalTok{k}
            \ControlFlowTok{if}\NormalTok{ (m }\OperatorTok{\textless{}}\NormalTok{ N):}
\NormalTok{                y[n] }\OperatorTok{+=}\NormalTok{ x[m]}
    \ControlFlowTok{return}\NormalTok{ y}
\end{Highlighting}
\end{Shaded}

Below we can see the plot for OSR=2, the right most plot is the
oversampled version.

The noise has all frequencies, and it's the high frequency components
that start to cancel each other. An average filter (sometimes called a
sinc filter due to the shape in the frequency domain) will have zeros at
\(\pm fs/2\) where the noise power tends towards zero.


{
\centering
\aicfig{media/l6_osr_2_tikz.png}
}

\emph{Figure 22: FFTs from continuous value to 10-bit quantized to
oversampled with OSR=2 (right), where the averaging filter nulls the
noise towards half the sample rate}

The low frequency components will add, and we can notice how the noise
power increases close to the zero frequency (middle of the x-axis).

For an OSR of 4 we can count four dips in the noise floor, although
there are really three nulls: a length-4 moving average has zeros at
\(f_s/4\), \(f_s/2\) and \(3f_s/4\), and on this two-sided plot the
\(f_s/2\) null is split across both edges so you see half of it twice.
For OSR=2 there is one null, at \(f_s/2\), and it shows up only as the
two edges.


{
\centering
\aicfig{media/l6_osr_4_tikz.png}
}

\emph{Figure 23: The same FFTs with OSR=4 (right), where the
moving-average filter puts three nulls in the noise floor, seen as four
dips on this two-sided plot, and the noise power increases close to zero
frequency}

The code for the plots is
\href{https://github.com/wulffern/aic2026/blob/main/ex/osr.py}{osr.py}.
I would encourage you to play a bit with the code, and make sure you
understand oversampling. If you would rather drag a slider than edit a
file, the
\href{https://wulffern.github.io/aic2026/assets/examples/oversampling.html}{interactive
version} plots the measured in-band SNR against OSR next to the ideal 3
dB per octave.

\section{Noise Shaping}\label{noise-shaping}\index{noise shaping@Noise shaping}

Look at the OSR=4 plot above, and look at it honestly, because it does
not show what you might hope. Near zero frequency the averaged floor is
not lower than the unaveraged one --- measured from the script it is
about 0.8 dB \emph{higher} for OSR=4, which is the same thing the
previous paragraph said when it noted that the low-frequency components
add. The averaging filter has done almost nothing to the noise in the
band we care about, because that noise is in its passband.

So where is the 6 dB that the theory promised? It is there, but it comes
from \emph{restricting the band}, not from the filter. Integrate the
same unaveraged spectrum over \(\vert f\vert < f_s/8\) instead of the
whole \(f_s/2\) and the signal-to-noise ratio improves by close to
\(10\log(4)\), exactly as the derivation said. The filter's job is not
to create that improvement but to make it safe to collect: it removes
the out-of-band noise that would otherwise fold back on top of the
signal when you decimate. These plots never decimate, so the payoff is
invisible in them. That is a limitation of the figure, not of
oversampling.

What the plots do show clearly is the ceiling. Even with the averaging,
the noise level of the discrete-time continuous-value plot is much lower
than anything the quantized paths achieve.

What if we could do something, add some circuitry, before the
quantization such that the quantization noise was reduced?

That's what noise shaping is all about. Adding circuits such that we can
``shape'' the quantization noise. We can't make the quantization noise
disappear, or indeed reduce the total noise power of the quantization
noise, but we can reduce the quantization noise power for a certain
frequency band.

But what circuitry can we add?

\subsection{The magic of feedback}\label{the-magic-of-feedback}\index{the magic of feedback@The magic of feedback}

A generalized feedback system is shown below, it could be a regulator, a
unity-gain buffer, or something else.

The output \(V_o\) is subtracted from the input \(V_i\), and the error
\(V_x\) is shaped by a filter \(H(s)\).

If we make \(H(s)\) infinite, then \(V_o = V_i\). If you've never seen
such a circuit, you might ask ``Why would we do this? Could we not just
use \(V_i\) directly?''. There are many reasons for using a circuit like
this, let me explain one instance.

Imagine we have a VDD of 1.8 V, and we want to make a 0.9 V voltage for
a CPU. The CPU can consume up to 10 mA. One way to make a divide by two
circuit is with two equal resistors connected between VDD and ground. We
don't want the resistive divider to consume a large current, so let's
choose 1 MOhm resistors. The current in the resistor divider would then
be about 1 \(\mu\)A. We can't connect the CPU directly to the resistor
divider, the CPU can draw 10 mA. As such, we need a copy of the voltage
at the mid-point of the resistor divider that can drive 10 mA.

Do you see now why a circuit like the one below is useful? If not, you
should really come talk to me so I can help you understand.


{
\centering
\aicfig{media/l4_sdloop_tikz.png}
}

\emph{Figure 24: A generalized feedback system where the error between
input and output is shaped by a filter H(s), and the output equals the
input when H(s) is infinite}

\subsection{Sigma-delta principle}\label{sigma-delta-principle}\index{sigma-delta principle@Sigma-delta principle}

Let's modify the feedback circuit into the one below. I've added an ADC
and a DAC to the feedback loop, and the \(D_o\) is now the output we're
interested in. The equation for the loop would be

\[
D_o = adc\left[H(s)\left(V_i - dac(D_o)\right)\right]
\]

But how can we now calculate the transfer function \(\frac{D_o}{V_i}\)?
Both \(adc\) and \(dac\) could be non-linear functions, so we can't
disentangle the equation. Let's make assumptions.


{
\centering
\aicfig{media/l4_sd_tikz.png}
}

\emph{Figure 25: The sigma-delta principle: a feedback loop with a
filter H(s), an ADC (quantizer) and a DAC in the feedback path, with
digital output Do}

\subsubsection{The DAC assumption}\label{the-dac-assumption}

\textbf{Assumption 1:} the \(dac\) is linear, such that
\(V_o = dac(D_o) = A  D_o + B\), where \(A\) and \(B\) are scalar
values.

The DAC must be linear, otherwise our noise-shaping ADC will not work.

One way to force linearity is to use a 1-bit DAC, which has only two
points, so should be linear. For example \[ V_o = A \times D_o\], where
\(D_o \in \{0,1\}\). Even a 1-bit DAC could be non-linear if \(A\) is
time-variant, so \(V_o[n] = A(t)\times D_o[n]\), this could happen if
the reference voltage for the DAC changed with time.

I've made a couple noise shaping ADCs, and in the first one I made I
screwed up the DAC. It turned out that the DAC current had a signal
dependent component which lead to a non-linear behavior.

\subsubsection{The ADC assumption}\label{the-adc-assumption}

\textbf{Assumption 2:} the \(adc\) can be modeled as a linear function
\(D_o = adc(x) = x + e\), where e is \textbf{white noise source}

We've talked about this, the \(e\) is not white, especially for low-bit
ADCs, so we usually have to add noise. Sometimes it's sufficient with
thermal noise, but often it's necessary to add a random, or
pseudo-random noise source at the input of the ADC.

\subsubsection{The modified equation}\label{the-modified-equation}

With the assumptions we can change the equation into

\[
D_o = adc\left[H(s)\left(V_i - dac(D_o)\right)\right] = H(s)\left( V_i - A D_o\right) + e
\]

In noise-shaping texts it's common to write the above equation as

\[
y = H(s)(u - y) + e
\]

or in the sample domain

\[ y[n] = e[n] + h*(u[n] - y[n])\]

which could be drawn in a signal flow graph as below.


{
\centering
\aicfig{media/l6_sdadc_tikz.png}
}

\emph{Figure 26: Signal flow graph of the noise-shaping loop: the
difference between input u{[}n{]} and output y{[}n{]} is filtered by
H(z) and the quantization noise e{[}n{]} is added at the quantizer}

in the Z-domain the equation would turn into

\[ Y(z) = E(z) + H(z)\left[U(z) - Y(z)\right]\]

The whole point of this exercise was to somehow shape the quantization
noise, and we're almost at the point, but to show how it works we need
to look at the transfer function for the signal \(U\) and for the noise
\(E\).

\subsection{Signal transfer function}\label{signal-transfer-function}\index{signal transfer function@Signal transfer function}

Assume U and E are uncorrelated, and E is zero

\[Y = HU - HY \]

\[ STF = \frac{Y}{U} = \frac{H}{1 + H} = \frac{1}{1 + \frac{1}{H}}\]

Imagine what will happen if H is infinite. Then the signal transfer
function (STF) is 1, and the output \(Y\) is equal to our input \(U\).
That's exactly what we wanted from the feedback circuit.

\subsection{Noise transfer function}\label{noise-transfer-function}\index{noise transfer function@Noise transfer function}

Assume U is zero

\[ Y = E - HY \rightarrow NTF = \frac{1}{1 + H}\]

Imagine again what happens when H is infinite. In this case the
noise-transfer function becomes zero. In other words, there is no added
noise.

\subsection{Combined transfer
function}\label{combined-transfer-function}

In the combined transfer function below, if we make \(H(z)\) infinite,
then \(Y = U\) and there is \textbf{no added quantization noise}. I
don't know how to make \(H(z)\) infinite everywhere, so we have to
choose at what frequencies it's ``infinite''.

\[Y(z) = STF(z) U(z) + NTF(z) E(z)\]

There are a large set of different \(H(z)\) and I'm sure engineers will
invent new ones. We usually classify the filters based on the number of
zeros in the NTF, for example, first-order (one zero), second order (two
zeros) etc. There are books written about sigma-delta modulators, and I
would encourage you to read those to get a deeper understanding. I would
start with \href{https://ieeexplore.ieee.org/book/5273726}{Delta-Sigma
Data Converters: Theory, Design, and Simulation}.

\section{First-Order Noise-Shaping}\label{first-order-noise-shaping}\index{first-order noise-shaping@First-order noise-shaping}

We want an infinite \(H(z)\). One way to get an infinite function is an
accumulator, for example

\[ y[n+1] = x[n] + y[n]\]

or in the Z-domain

\[ zY = X + Y \rightarrow Y(z-1) = X\]

which has the transfer function

\[H(z) = \frac{1}{z-1}\]

The signal transfer function is

\[STF = \frac{1/(z-1)}{1 + 1/(z-1)} = \frac{1}{z} = z^{-1}\]

and the noise transfer function

\[NTF = \frac{1}{1 + 1/(z-1)} = \frac{z-1}{z} = 1 - z^{-1}\]

In order calculate the Signal to Quantization Noise Ratio we need to
have an expression for how the NTF above filters the quantization noise.

In the book they replace the \(z\) by the Laplace variable and then walk
out onto the imaginary axis, \(s = j\omega\), which is where a transfer
function turns into a frequency response.

\[z = e^{sT} \;\underset{s=j\omega}{\longrightarrow}\;  e^{j\omega T} = e^{j2 \pi f/f_s}\]

inserted into the NTF we get the function below. The three lines are one
trick applied once, so it is worth knowing what you are looking for
before reading them. We want \(1 - e^{-j\theta}\) turned into something
whose magnitude we can read off, and the only identity available is
\(\sin\theta = (e^{j\theta}-e^{-j\theta})/2j\). So we factor out half
the exponent, \(e^{-j\pi f/f_s}\), which leaves a difference of two
conjugate exponentials in the bracket --- exactly the shape of that
identity --- and then divide and multiply by \(2j\) to make it literally
a sine. Everything pulled out has magnitude 1, so it contributes phase
and nothing else.

\[\begin{aligned}
NTF(f) &= 1- e^{-j2 \pi f/f_s} \\
       &= \frac{e^{j \pi f/f_s} -e^{-j \pi f/f_s}}{2j}\times 2j \times e^{-j\pi f/f_s} \\
       &= \sin\left(\frac{\pi f}{f_s}\right) \times 2j \times e^{-j \pi f/f_s}
\end{aligned}\]

The arithmetic magic is really to extract the
\(2j \times e^{-j \pi f/f_s}\) from the first expression such that the
initial part can be translated into a sinusoid.

When we take the absolute value to figure out how the NTF changes with
frequency the complex parts disappears (equal to 1)

\[\vert NTF(f)\vert = \left\vert 2 \sin\left(\frac{\pi f}{f_s}\right)\right\vert \]

The signal power for a sinusoid is

\[ P_s = A^2/2\]

The in-band noise power for the shaped quantization noise is

\[ P_n = \int_{-f_0}^{f_0} \frac{\Delta^2}{12}\frac{1}{f_s}\left[2 \sin\left(\frac{\pi f}{f_s}\right)\right]^2 df\]

The integral is not actually tedious, and it is worth doing once because
it explains both of the odd-looking constants in the answer. In band,
\(f\) is much smaller than \(f_s\), so
\(\sin(\pi f/f_s) \approx \pi f/f_s\) and the integrand becomes a
parabola:

\[ P_n \approx \frac{\Delta^2}{12}\frac{1}{f_s}\frac{4\pi^2}{f_s^2}\int_{-f_0}^{f_0} f^2 df = \frac{\Delta^2}{12}\frac{4\pi^2}{f_s^3}\frac{2f_0^3}{3} \]

That \(f_0^3\) is where the whole advantage comes from. Ordinary
oversampling had the in-band noise falling as \(f_0\); here it falls as
\(f_0^3\), because the shaping makes the noise density itself
proportional to \(f^2\). Substituting \(OSR = f_s/2f_0\):

\[ P_n \approx \frac{\Delta^2}{12}\frac{\pi^2}{3}\frac{1}{OSR^3} \]

The \(OSR^3\) becomes the \(30\log(OSR)\) --- three times the
\(10\log(OSR)\) of plain oversampling --- and the \(\pi^2/3\) becomes
the penalty \(10\log(\pi^2/3) = 5.17\) dB. Take the ratio to
\(P_s = A^2/2\) and

\[SQNR = 6.02 B + 1.76 - 5.17 + 30 \log(OSR)\]

If we compare to pure oversampling, where the SQNR improves by
\(10 \log(OSR)\), a first order sigma-delta improves by
\(30 \log(OSR)\). That's a significant improvement.

\subsection{SQNR and ENOB}\label{sqnr-and-enob}\index{sqnr and enob@SQNR and ENOB}

Below is the signal-to-quantization noise ratio's for Nyquist up to
second order sigma-delta. Read them as one family. Every line starts
from the same \(6.02B + 1.76\), and each step down the list buys a
steeper dependence on OSR at the cost of a larger fixed penalty.

The second-order line follows from repeating the derivation above with
\(NTF = (1-z^{-1})^2\), so the noise density goes as \(f^4\) instead of
\(f^2\). The integral then produces \(OSR^5\), hence \(50\log(OSR)\),
and \(\pi^4/5\) in place of \(\pi^2/3\), hence
\(10\log(\pi^4/5) = 12.9\) dB. The pattern continues: an \(L\)'th order
modulator gives \((20L+10)\log(OSR)\) and a penalty of
\(10\log(\pi^{2L}/(2L+1))\).

The penalties are real, and at low OSR they can outweigh the steeper
slope. Setting the first- and second-order expressions equal gives a
crossover at \(OSR \approx 2.4\): below that, second-order shaping is
\emph{worse} than first-order, because the extra 7.7 dB of penalty has
not yet been earned back by the extra \(20\log(OSR)\). Higher order is
not automatically better, it is better \emph{eventually}, and where
``eventually'' starts is a number you can compute before committing to
an architecture.

\[SQNR_{nyquist} \approx 6.02B + 1.76 \]

\[SQNR_{oversample} \approx 6.02B + 1.76 + 10 \log(OSR)\]

\[SQNR_{\Sigma\Delta 1} \approx 6.02 B + 1.76 - 5.17 + 30 \log(OSR)\]

\[SQNR_{\Sigma\Delta 2} \approx 6.02 B + 1.76 - 12.9 + 50 \log(OSR)\]

We could compute an effective number of bits, as shown below.

\[ ENOB = (SQNR - 1.76)/6.02 \]

Assume 1-bit quantizer, what would be the maximum ENOB?

\begin{longtable}[]{@{}cccc@{}}
\toprule\noalign{}
OSR & Oversampling & First-order & Second-order \\
\midrule\noalign{}
\endhead
\bottomrule\noalign{}
\endlastfoot
4 & 2.0 & 3.1 & 3.9 \\
64 & 4.0 & 9.1 & 13.9 \\
1024 & 6.0 & 15.1 & 23.9 \\
\end{longtable}

\emph{Table: ENOB of a 1-bit quantizer (\(B = 1\)), from the three
expressions above.}

Read down a column and you see what oversampling buys you; read across a
row and you see what an architecture buys you. Plain oversampling is
hopeless: a thousand-fold increase in sample rate turns one bit into
six. First-order shaping turns the same thousand-fold into fifteen, and
second-order into twenty-four. That is the entire argument for building
the loop, and it is why a 1-bit modulator followed by a decimation
filter is a sensible way to make a 20-bit converter, while a 1-bit
oversampled ADC is not a way to make anything.

The second-order column also shows the crossover discussed above. At
\(OSR=4\) second-order leads first-order by less than a bit, because the
extra 7.7 dB penalty has barely been paid off; by \(OSR=1024\) it leads
by nine.

\section{Examples}\label{examples}\index{examples@Examples}

\subsection{Python noise-shaping}\label{python-noise-shaping}\index{python noise-shaping@Python noise-shaping}

I want to demystify noise-shaping modulators. I think one way to do that
is to show some code. You can find the code at
\href{https://github.com/wulffern/aic2026/blob/main/ex/sd_1st.py}{sd\_1st.py},
and an
\href{https://wulffern.github.io/aic2026/assets/examples/sigma-delta.html}{interactive
version} that runs the same loop in your browser and decodes the
bitstream back into the input.

Below we can see an excerpt. Again pretty stupid code, and I'm sure it's
possible to make a faster version (for loops in python are notoriously
slow).

For each sample in the input vector \(u\) I compute the input to the
quantizer \(x\), which is the sum of the previous input to the quantizer
and the difference between the current input and the previous output
\(y_{sd}\).

The quantizer generates the next \(y_{sd}\) and I have the option to add
dither.

Three details in the full script are not in this excerpt and matter if
you want to reproduce the plots. The input is scaled to 0.7 of full
scale, because a true 1-bit loop overloads if you drive it all the way
--- the integrator runs away and the output degenerates into a slow
square wave. The first two samples are thrown away, because the
integrator starts at zero and needs a moment to settle. And the dither
amplitude is a quarter of a quantizer step, which is enough to break up
idle tones without swamping the signal.

One inconsistency is worth flagging rather than hiding. This
\texttt{quantize} uses a step of \(2/(2^B-1)\) for more than one bit, a
mid-tread converter whose levels include zero, while the \texttt{adc}
function earlier in this chapter uses \(\Delta = 2/2^B\), a mid-rise
converter with no zero level. Both are real converters and both appear
in real silicon; they differ by half a step and by whether a zero input
produces a zero output. The \(\Delta^2/12\) result holds for either.
Just do not mix the two definitions when you are comparing measured
numbers against theory, which is exactly the sort of half-LSB
discrepancy that costs an afternoon.

\begin{Shaded}
\begin{Highlighting}[]
\KeywordTok{def}\NormalTok{ quantize(v,bits):}
    \CommentTok{\#{-} 2**bits levels reaching +/{-}1, so bits=1 is}
    \CommentTok{\#{-} a genuine two{-}level quantizer}
\NormalTok{    levels }\OperatorTok{=} \DecValTok{2}\OperatorTok{**}\NormalTok{bits}
    \ControlFlowTok{if}\NormalTok{(levels }\OperatorTok{==} \DecValTok{2}\NormalTok{):}
        \ControlFlowTok{return} \FloatTok{1.0} \ControlFlowTok{if}\NormalTok{ v }\OperatorTok{\textgreater{}=} \DecValTok{0} \ControlFlowTok{else} \OperatorTok{{-}}\FloatTok{1.0}
\NormalTok{    step }\OperatorTok{=} \DecValTok{2}\OperatorTok{/}\NormalTok{(levels}\OperatorTok{{-}}\DecValTok{1}\NormalTok{)}
    \ControlFlowTok{return} \BuiltInTok{float}\NormalTok{(np.clip(np.}\BuiltInTok{round}\NormalTok{(v}\OperatorTok{/}\NormalTok{step)}\OperatorTok{*}\NormalTok{step,}\OperatorTok{{-}}\DecValTok{1}\NormalTok{,}\DecValTok{1}\NormalTok{))}

\CommentTok{\# u is discrete time, continuous value input}
\NormalTok{M }\OperatorTok{=} \BuiltInTok{len}\NormalTok{(u)}
\NormalTok{y\_sd }\OperatorTok{=}\NormalTok{ np.zeros(M)}
\NormalTok{x }\OperatorTok{=}\NormalTok{ np.zeros(M)}
\ControlFlowTok{for}\NormalTok{ n }\KeywordTok{in} \BuiltInTok{range}\NormalTok{(}\DecValTok{1}\NormalTok{,M):}
\NormalTok{    x[n] }\OperatorTok{=}\NormalTok{ x[n}\OperatorTok{{-}}\DecValTok{1}\NormalTok{] }\OperatorTok{+}\NormalTok{ (u[n]}\OperatorTok{{-}}\NormalTok{y\_sd[n}\OperatorTok{{-}}\DecValTok{1}\NormalTok{])}
\NormalTok{    y\_sd[n] }\OperatorTok{=}\NormalTok{ quantize(x[n]}
        \OperatorTok{+}\NormalTok{ dither}\OperatorTok{*}\NormalTok{np.random.randn()}\OperatorTok{/}\NormalTok{(}\DecValTok{4}\OperatorTok{*}\DecValTok{2}\OperatorTok{**}\NormalTok{bits),bits)}
\end{Highlighting}
\end{Shaded}

The right-most plot is the one with noise-shaping. We can observe that
the noise seems to tend towards zero at zero frequency, as we would
expect. The accumulator above would have an infinite gain at infinite
time (it's the sum of all previous values), as such, the NTF goes
towards zero at 0 frequency.

If we look at the noise we can also see the non-white quantization
noise, which will degrade our performance. I hope by now, you've grown
tired of me harping on the point that \textbf{quantization noise is not
white}


{
\centering
\aicfig{media/l6_sd_d0_b1_tikz.png}
}

\emph{Figure 27: First-order sigma-delta modulator with 1-bit quantizer
and no dither, where the noise-shaped spectrum (right) tends towards
zero at zero frequency but contains distinct tones}

In the figure below I've turned on dither, and we can see how the noise
looks ``better'', which I know is not a qualitative statement, but ask
anyone that's done 1-bit quantizers. It's important to have enough
random noise.

Be clear about what has been bought and what has been paid, though,
because the y-axis of the plot tells both halves of the story. The
distinct tones are gone and the floor is smooth, which is what
``better'' means here: a smooth floor is predictable, it does not move
when the input does, and it will not land a spur in the middle of your
signal band on a Tuesday. But the floor is also \emph{higher}, and the
notch at zero frequency is shallower. Running the loop with and without
dither, the in-band signal-to-noise ratio drops by a decibel or two.
Dither is not free noise reduction; it converts a small amount of
signal-to-noise ratio into a large amount of predictability. That is
usually a good trade for a 1-bit quantizer with a slow input, where the
alternative is idle tones parked wherever the input DC level happens to
put them, and a bad trade for a many-bit quantizer that was never going
to produce tones in the first place.


{
\centering
\aicfig{media/l6_sd_d1_b1_tikz.png}
}

\emph{Figure 28: The same first-order 1-bit sigma-delta modulator with
dither enabled, where the noise-shaped spectrum (right) is smoother and
more noise-like}

In papers it's common to use a logarithmic frequency axis, as shown
below. In the plot I only show the positive frequencies of the FFT. From
the shape of the quantization noise we can also see the first order
behavior, as a straight 20 dB per decade rise, which is much easier to
read off a log axis than off the linear ones above.

Two things have changed from the previous two figures, and it is worth
saying so rather than letting you wonder. The quantizer here is 5-bit,
not 1-bit, so the floor is far lower and there are no idle tones to
speak of. And the y-axis is a magnitude spectrum in dB, not a power
spectral density: nothing has been normalised to the bin width, so the
absolute level is only meaningful relative to the tone.


{
\centering
\aicfig{media/l6_sdlog_d1_b5_tikz.png}
}

\emph{Figure 29: Magnitude spectrum of the output of a 5-bit first-order
sigma-delta modulator on a logarithmic frequency axis, showing the 20
dB/decade shaping of the quantization noise}

\subsection{The wonderful world of SD
modulators}\label{the-wonderful-world-of-sd-modulators}

\subsubsection{Open-Loop Sigma-Delta}\label{open-loop-sigma-delta}

On my Ph.D I did some work on

Resonators in Open-Loop Sigma-Delta Modulators
\citeproc{ref-wulff09}{{[}4{]}}

which was a pure theoretical work. The idea was to use modulo
integrators (local control of integrator output swing) in front of large
latency multi-bit quantizers to achieve a high SNR.

The plot below shows a fifth order NTF where there are two complex
conjugate zero \emph{pairs}, and a zero at zero frequency. With a higher
order filter one can use a lower OSR, and still achieve high ENOB.


{
\centering
\aicfig{media/l06_osd21.png}
}

\emph{Figure 30: Output spectrum of an open-loop sigma-delta modulator
with a fifth-order NTF (two complex conjugate zero pairs and a zero at
DC), reaching 13.8 bit ENOB and 84.9 dB SNDR}

\subsubsection{Noise Shaped SAR}\label{noise-shaped-sar}

One of my Ph.d students made a

A 68 dB SNDR Compiled Noise-Shaping SAR ADC With On-Chip CDAC
Calibration \citeproc{ref-garvik19}{{[}5{]}}

In a SAR ADC, once the bit-cycling is complete, the analog value on the
capacitors is the actual quantization error. That error can be fed to a
loop filter, H(z), and amplified in the next conversion, accordingly a
combination of SAR and noise-shaping.

In the paper the SD modulator was also used to calibrate the
non-linearity in the CDAC, as the MSB capacitor won't be exactly N times
larger than the smallest capacitor.


{
\centering
\aicfig{media/l6_harald_arch.png}
}

\emph{Figure 31: Architecture of the noise-shaping SAR ADC: capacitive
DAC with multiplexers, loop filter H(z), integrating comparator, SAR
logic, calibration logic and code correction}

The loop filter was a switched cap loop filter, and we can see the NTF
below. The first OTA made use of chopping to reduce the offset.


{
\centering
\aicfig{media/l6_fig_harald_circuit.png}
}

\emph{Figure 32: The switched-capacitor loop filter with two OTAs (the
first one chopped), the clock phases relative to the SAR activity, and
the resulting NTF with -27.8 dB in-band suppression}

\subsubsection{Control-Bounded ADCs}\label{control-bounded-adcs}

One of my current Ph.D students is working on an even more advanced type
of sigma-delta ADC. Actually, it's more a super-set of SD ADCs called
control-bounded ADCs.

\href{https://ntnuopen.ntnu.no/ntnu-xmlui/handle/11250/2824253}{Design
Considerations for a Low-Power Control-Bounded A/D Converter}

A block diagram of a Leapfrog ADC version of a control-bounded ADC is
shown below.

Here we're walking into advanced maths territory, but to simplify, I
think it's correct to say that a control-bounded ADC seeks to control
the local analog state, \(x_n(t)\) such that no voltage is saturated.
The digital control signals \(s_n(t)\) are used to infer the state of
the input \(u(t)\) using a form of
\href{https://en.wikipedia.org/wiki/Bayesian_statistics}{Bayesian
Statistics}.


{
\centering
\aicfig{media/l6_fredrik_arch.png}
}

\emph{Figure 33: Block diagram of the Leapfrog control-bounded ADC: a
chain of continuous-time integrators with local digital control loops
s(t) that keep the analog states x(t) bounded}

Below we can see a power spectral density plot of the ADC, and we can
observe how the quantization noise is shaped. I think it's a third order
NTF with a zero at zero frequency and a complex conjugate zero pair, a
notch, at 8 MHzish.


{
\centering
\aicfig{media/l6_fredrik_psd.png}
}

\emph{Figure 34: Power spectral density of the control-bounded ADC's
estimated input together with the NTF, a third-order shaping with a
notch around 8 MHz}

\subsubsection{Complex Sigma-Delta}\label{complex-sigma-delta}

There are cool sigma-delta modulators with crazy configurations and that
may look like an exercise in ``Let's make something complex'', however,
most of them have a reasonable application. One example is the one below
for radio receivers

\href{https://ieeexplore.ieee.org/stamp/stamp.jsp?tp=&arnumber=4381437}{A
56 mW Continuous-Time Quadrature Cascaded Sigma-Delta Modulator With 77
dB DR in a Near Zero-IF 20 MHz Band}


{
\centering
\aicfig{media/qt_sd_tikz.png}
}

\emph{Figure 35: One stage of the quadrature modulator, with the second
cascaded stage and the differential wiring left out. Inspired by Breems
et al. \citeproc{ref-breems07}{{[}6{]}}}

The cross-coupling in red is the whole of what makes it quadrature. Take
it away and there are two independent real modulators, each with a noise
transfer function that is symmetric about zero frequency, because a real
circuit cannot tell \(+f\) from \(-f\). Put it back and the pair
responds to a rotating phasor rather than an oscillating one, so the
notch can sit on one side of zero and not the other. That is exactly
what a near zero-IF receiver wants: the wanted channel is a little above
zero, its image a little below, and only a complex modulator can quieten
one without quietening both.

\subsubsection{My first Sigma-Delta}\label{my-first-sigma-delta}

The first sigma-delta modulator I made in ``real-life'' was similar to
the one shown below.

The input voltage is translated into a current, and the current is
integrated on capacitor \(C\). The \(R_{offset}\) is to change the
mid-level voltage, while \(R_{ref}\) is the 1-bit feedback DAC. The
comparator is the quantizer. When the clock strikes the comparator
compares the \(V_o\) and \(V_{ref}/2\) and outputs a 1-bit digital
output \(D\)

The complete ADC is operated in a ``incremental mode'', which is a fancy
way of saying

\begin{quote}
Reset your sigma-delta modulator, run the sigma delta modulator for a
fixed number of cycles (i.e 1024), and count the number of ones at \(D\)
\end{quote}

The effect of an ``incremental mode'' is to combine the modulator and a
output filter so the ADC appears to be a slow Nyquist ADC.

For more information, ask me, or see the patent at
\href{https://patents.google.com/patent/US8947280B2/en?inventor=carsten+wulff&oq=carsten+wulff}{Analogue-to-digital
converter}


{
\centering
\aicfig{media/l6_patent.png}
}

\emph{Figure 36: Incremental first-order sigma-delta ADC from the
patent: input and reference resistors into an OTA integrating on C, a
clocked comparator as quantizer, and a counter as output filter}

\section{Summary}\label{summary}

The one-page version of this chapter:

\begin{itemize}
\tightlist
\item
  One number for an ADC: the figure of merit - Walden for speed-limited,
  Schreier for noise-limited designs
\item
  Ideal quantization gives SQNR = 6.02 N + 1.76 dB, and the white-noise
  model of it holds only for busy inputs
\item
  Oversampling spreads the same noise power over more bandwidth: 3 dB
  (half a bit) per octave
\item
  Feedback around the quantizer shapes the noise away from the band:
  first-order sigma-delta buys 9 dB per octave, and higher order buys
  more
\item
  The decimation filter is where the promised resolution is actually
  cashed out
\item
  A compiled ADC is a netlist, an object file and a rule file - portable
  across processes in weeks, and good enough for JSSC
\end{itemize}

\section{Would you like to know
more?}\label{would-you-like-to-know-more}

The design of sigma-delta modulation analog-to-digital converters
\citeproc{ref-boser88}{{[}7{]}}

Delta-sigma modulation in fractional-N frequency synthesis
\citeproc{ref-riley93}{{[}8{]}}

A CMOS Temperature Sensor With a Voltage-Calibrated Inaccuracy of ± 0.15
C (3sigma) From -55 Cto 125 C \citeproc{ref-souri13}{{[}9{]}}

A 20-mW 640-MHz CMOS Continuous-Time Sigma-Delta ADC With 20-MHz Signal
Bandwidth, 80-dB Dynamic Range and 12-bit ENOB
\citeproc{ref-mitteregger06a}{{[}10{]}}

A Micro-Power Two-Step Incremental Analog-to-Digital Converter
\citeproc{ref-chen15}{{[}11{]}}

\phantomsection\label{refs}
\begin{CSLReferences}{0}{0}
\bibitem[\citeproctext]{ref-walden99}
\CSLLeftMargin{{[}1{]} }%
\CSLRightInline{R. H. Walden, {``Analog to digital converter survey and
analysis,''} \emph{IEEE Journal on Selected Areas in Communications},
vol. 17, no. 4, pp. 539--550, 1999, doi:
\href{https://doi.org/10.1109/49.761034}{10.1109/49.761034}. }

\bibitem[\citeproctext]{ref-wulff17}
\CSLLeftMargin{{[}2{]} }%
\CSLRightInline{C. Wulff and T. Ytterdal, {``A compiled 9-bit 20-MS/s
3.5-fJ/conv.step SAR ADC in 28-nm FDSOI for bluetooth low energy
receivers,''} \emph{IEEE Journal of Solid-State Circuits}, vol. 52, no.
7, pp. 1915--1926, 2017, doi:
\href{https://doi.org/10.1109/JSSC.2017.2685463}{10.1109/JSSC.2017.2685463}.
}

\bibitem[\citeproctext]{ref-blachman85a}
\CSLLeftMargin{{[}3{]} }%
\CSLRightInline{N. Blachman, {``The intermodulation and distortion due
to quantization of sinusoids,''} \emph{IEEE Transactions on Acoustics,
Speech, and Signal Processing}, vol. 33, no. 6, pp. 1417--1426, 1985,
doi:
\href{https://doi.org/10.1109/TASSP.1985.1164729}{10.1109/TASSP.1985.1164729}.
}

\bibitem[\citeproctext]{ref-wulff09}
\CSLLeftMargin{{[}4{]} }%
\CSLRightInline{C. Wulff and T. Ytterdal, {``Resonators in open-loop
sigma--delta modulators,''} \emph{IEEE Transactions on Circuits and
Systems I: Regular Papers}, vol. 56, no. 10, pp. 2159--2172, 2009, doi:
\href{https://doi.org/10.1109/TCSI.2009.2015211}{10.1109/TCSI.2009.2015211}.
}

\bibitem[\citeproctext]{ref-garvik19}
\CSLLeftMargin{{[}5{]} }%
\CSLRightInline{H. Garvik, C. Wulff, and T. Ytterdal, {``A 68 dB SNDR
compiled noise-shaping SAR ADC with on-chip CDAC calibration,''} in
\emph{2019 IEEE asian solid-state circuits conference (a-SSCC)}, 2019,
pp. 193--194, doi:
\href{https://doi.org/10.1109/A-SSCC47793.2019.9056925}{10.1109/A-SSCC47793.2019.9056925}.
}

\bibitem[\citeproctext]{ref-breems07}
\CSLLeftMargin{{[}6{]} }%
\CSLRightInline{L. J. Breems, R. Rutten, R. H. M. van Veldhoven, and G.
van der Weide, {``A 56 mW continuous-time quadrature cascaded
\(\Sigma\Delta\) modulator with 77 dB DR in a near zero-IF 20 MHz
band,''} \emph{IEEE Journal of Solid-State Circuits}, vol. 42, no. 12,
pp. 2696--2705, 2007, doi:
\href{https://doi.org/10.1109/JSSC.2007.908765}{10.1109/JSSC.2007.908765}.
}

\bibitem[\citeproctext]{ref-boser88}
\CSLLeftMargin{{[}7{]} }%
\CSLRightInline{B. Boser and B. Wooley, {``The design of sigma-delta
modulation analog-to-digital converters,''} \emph{IEEE Journal of
Solid-State Circuits}, vol. 23, no. 6, pp. 1298--1308, 1988, doi:
\href{https://doi.org/10.1109/4.90025}{10.1109/4.90025}. }

\bibitem[\citeproctext]{ref-riley93}
\CSLLeftMargin{{[}8{]} }%
\CSLRightInline{T. Riley, M. Copeland, and T. Kwasniewski,
{``Delta-sigma modulation in fractional-n frequency synthesis,''}
\emph{IEEE Journal of Solid-State Circuits}, vol. 28, no. 5, pp.
553--559, 1993, doi:
\href{https://doi.org/10.1109/4.229400}{10.1109/4.229400}. }

\bibitem[\citeproctext]{ref-souri13}
\CSLLeftMargin{{[}9{]} }%
\CSLRightInline{K. Souri, Y. Chae, and K. A. A. Makinwa, {``A CMOS
temperature sensor with a voltage-calibrated inaccuracy of \(\pm\)
0.15\(^{\circ}\) c (3\(\sigma\) ) from \(-\) 55\(^{\circ}\) c to
125\(^{\circ}\) c,''} \emph{IEEE Journal of Solid-State Circuits}, vol.
48, no. 1, pp. 292--301, 2013, doi:
\href{https://doi.org/10.1109/JSSC.2012.2214831}{10.1109/JSSC.2012.2214831}.
}

\bibitem[\citeproctext]{ref-mitteregger06a}
\CSLLeftMargin{{[}10{]} }%
\CSLRightInline{G. Mitteregger, C. Ebner, S. Mechnig, T. Blon, C.
Holuigue, and E. Romani, {``A 20-mW 640-MHz CMOS continuous-time
\(\Sigma\Delta\) ADC with 20-MHz signal bandwidth, 80-dB dynamic range
and 12-bit ENOB,''} \emph{IEEE Journal of Solid-State Circuits}, vol.
41, no. 12, pp. 2641--2649, 2006, doi:
\href{https://doi.org/10.1109/JSSC.2006.884332}{10.1109/JSSC.2006.884332}.
}

\bibitem[\citeproctext]{ref-chen15}
\CSLLeftMargin{{[}11{]} }%
\CSLRightInline{C.-H. Chen, Y. Zhang, T. He, P. Y. Chiang, and G. C.
Temes, {``A micro-power two-step incremental analog-to-digital
converter,''} \emph{IEEE Journal of Solid-State Circuits}, vol. 50, no.
8, pp. 1796--1808, 2015, doi:
\href{https://doi.org/10.1109/JSSC.2015.2413842}{10.1109/JSSC.2015.2413842}.
}

\end{CSLReferences}
